\documentclass[10pt]{extarticle}
\usepackage[a4paper,margin=1.5cm]{geometry}
\usepackage[T1]{fontenc}
\usepackage{amsmath,amssymb}
\usepackage{graphicx}
\usepackage[export]{adjustbox}
\usepackage{multicol}
\usepackage{enumitem}
\usepackage{booktabs,array,tabularx}
\setlength{\columnsep}{0.7cm}
\setlength{\parindent}{0pt}
\setlength{\parskip}{3pt}
\pagestyle{empty}
\makeatletter
% figure registry (figdefs.tex is generated by the extract step)
\newcommand{\qfigdef}[4]{\@namedef{qf@file@#1}{#2}\@namedef{qf@w@#1}{#3}\@namedef{qf@wide@#1}{#4}}
\newcommand{\qf@check}[1]{\@ifundefined{qf@file@#1}{\PackageError{score2}{Unknown figure id #1}{}}{}}
\newcommand{\qf@img}[4][]{\includegraphics[width=#2\dimexpr\@nameuse{qf@w@#4}pt\relax,max width=#3,max height=0.45\textheight,#1]{\@nameuse{qf@file@#4}}}
% Full-width questions: finalize wraps every question that contains a wide figure in \qfullbegin...\qfullend,
% so the WHOLE question (stem, figures, options) is set across the page in reading order, then the two columns
% resume. The question's head - its number, stem and first figure (\qfullsplit marks the end of the head) -
% is kept on one page, so a number is never left at a page foot with its figure overleaf; the rest may break
% across pages normally, so no big gaps are left. Without a \qfullsplit, the whole question is kept together
% if it is under 45% of a page. Never write these yourself.
\newif\ifqf@full
\newif\ifqf@boxed
\newsavebox{\qf@box}
\newcommand{\qf@open}{\global\qf@fulltrue\par\global\qf@boxedtrue
  \setbox\qf@box\vbox\bgroup\hsize\textwidth\linewidth\textwidth}
\newcommand{\qf@emit}[1]{\par\egroup\global\qf@boxedfalse
  \ifdim\dimexpr\ht\qf@box+\dp\qf@box\relax<#1\textheight\noindent\usebox{\qf@box}\par\else\unvbox\qf@box\fi}
\newcommand{\qfullsplit}{\ifqf@boxed\qf@emit{0.85}\fi}
\newcommand{\qf@close}{\ifqf@boxed\qf@emit{0.45}\fi\global\qf@fullfalse}
\newcommand{\qfullbegin}{\end{multicols}\qf@open}
\newcommand{\qfullend}{\qf@close\begin{multicols}{2}}
\newcommand{\qfullstart}{\qf@open}
\newcommand{\qfulldone}{\qf@close}
% \qfig[scale]{id}: block figure, centred, directly where it is written (size set by the tool).
% Fallback if a wide figure is not inside a full-width question: break out of the columns in place - no floats.
\newcommand{\qfig}[2][1]{\qf@check{#2}%
  \ifqf@full\qf@blk{#1}{#2}\else
  \ifnum\@nameuse{qf@wide@#2}=1 \expandafter\@firstoftwo\else\expandafter\@secondoftwo\fi
  {\end{multicols}\begin{center}\qf@img{#1}{\textwidth}{#2}\end{center}\begin{multicols}{2}}%
  {\qf@blk{#1}{#2}}\fi}
\newcommand{\qf@blk}[2]{\par\smallskip{\centering\qf@img{#1}{\linewidth}{#2}\par}\smallskip}
% \qfigi[scale]{id}: inline figure (inside option lists / tables / minipages); never breaks columns.
\newcommand{\qfigi}[2][1]{\qf@check{#2}\qf@img[valign=c]{#1}{\linewidth}{#2}}
\makeatother
% structure macros - lint reads these to check completeness and figure placement
\newcommand{\papersec}[1]{\par\bigskip{\centering\Large\bfseries PAPER-#1\par}\medskip}
\newcommand{\partsec}[2]{\par\medskip\noindent{\large\bfseries PART-#1: #2}\par\smallskip}
\newcommand{\sect}[1]{\par\smallskip\noindent\textbf{#1}\par}
\newcommand{\qn}[1]{\par\smallskip\noindent\textbf{#1.}~}
\newcommand{\paraq}[2]{\par\smallskip\noindent\textbf{Paragraph for Questions #1 and #2}\par}
\newenvironment{qoptions}{\begin{itemize}[leftmargin=1.8em,itemsep=1pt,topsep=2pt,label={}]}{\end{itemize}}
\newcommand{\opt}[2]{\item[(#1)] #2}
% figure registry (generated by score2_tools.py finalize)
\qfigdef{p1_phy_s1_q04_1}{TOPS - Test 02 (28 Aug 2022)-img/p1_phy_s1_q04_1.png}{142.7}{0}
\qfigdef{p1_phy_s1_q05_1}{TOPS - Test 02 (28 Aug 2022)-img/p1_phy_s1_q05_1.png}{82.0}{0}
\qfigdef{p1_phy_s2_q07_1}{TOPS - Test 02 (28 Aug 2022)-img/p1_phy_s2_q07_1.png}{100.2}{0}
\qfigdef{p1_phy_s2_q08_1}{TOPS - Test 02 (28 Aug 2022)-img/p1_phy_s2_q08_1.png}{198.1}{0}
\qfigdef{p1_phy_s2_q10_1}{TOPS - Test 02 (28 Aug 2022)-img/p1_phy_s2_q10_1.png}{108.3}{0}
\qfigdef{p1_phy_s2_q11_1}{TOPS - Test 02 (28 Aug 2022)-img/p1_phy_s2_q11_1.png}{119.0}{0}
\qfigdef{p1_phy_s2_q12_1}{TOPS - Test 02 (28 Aug 2022)-img/p1_phy_s2_q12_1.png}{152.0}{0}
\qfigdef{p1_phy_s3_q01_1}{TOPS - Test 02 (28 Aug 2022)-img/p1_phy_s3_q01_1.png}{137.4}{0}
\qfigdef{p1_phy_s3_q02_1}{TOPS - Test 02 (28 Aug 2022)-img/p1_phy_s3_q02_1.png}{111.2}{0}
\qfigdef{p1_phy_s3_q03_1}{TOPS - Test 02 (28 Aug 2022)-img/p1_phy_s3_q03_1.png}{66.1}{0}
\qfigdef{p1_phy_s3_q04_1}{TOPS - Test 02 (28 Aug 2022)-img/p1_phy_s3_q04_1.png}{134.5}{0}
\qfigdef{p1_phy_s3_q05_1}{TOPS - Test 02 (28 Aug 2022)-img/p1_phy_s3_q05_1.png}{187.3}{0}
\qfigdef{p1_phy_s3_q06_1}{TOPS - Test 02 (28 Aug 2022)-img/p1_phy_s3_q06_1.png}{108.3}{0}
\qfigdef{p1_chem_s3_q01_1}{TOPS - Test 02 (28 Aug 2022)-img/p1_chem_s3_q01_1.png}{115.3}{0}
\qfigdef{p2_phy_s1_q02_1}{TOPS - Test 02 (28 Aug 2022)-img/p2_phy_s1_q02_1.png}{127.6}{0}
\qfigdef{p2_phy_s1_q03_1}{TOPS - Test 02 (28 Aug 2022)-img/p2_phy_s1_q03_1.png}{108.3}{0}
\qfigdef{p2_phy_s1_q04_1}{TOPS - Test 02 (28 Aug 2022)-img/p2_phy_s1_q04_1.png}{159.2}{0}
\qfigdef{p2_phy_s1_q05_1}{TOPS - Test 02 (28 Aug 2022)-img/p2_phy_s1_q05_1.png}{81.1}{0}
\qfigdef{p2_phy_s1_q06_1}{TOPS - Test 02 (28 Aug 2022)-img/p2_phy_s1_q06_1.png}{97.7}{0}
\qfigdef{p2_phy_s1_q06_2}{TOPS - Test 02 (28 Aug 2022)-img/p2_phy_s1_q06_2.png}{103.2}{0}
\qfigdef{p2_phy_s2_q02_1}{TOPS - Test 02 (28 Aug 2022)-img/p2_phy_s2_q02_1.png}{100.2}{0}
\qfigdef{p2_phy_s2_q03_1}{TOPS - Test 02 (28 Aug 2022)-img/p2_phy_s2_q03_1.png}{90.3}{0}
\qfigdef{p2_phy_s2_q04_1}{TOPS - Test 02 (28 Aug 2022)-img/p2_phy_s2_q04_1.png}{83.0}{0}
\qfigdef{p2_phy_s2_q05_1}{TOPS - Test 02 (28 Aug 2022)-img/p2_phy_s2_q05_1.png}{194.5}{0}
\qfigdef{p2_phy_s3_q04_1}{TOPS - Test 02 (28 Aug 2022)-img/p2_phy_s3_q04_1.png}{68.9}{0}
\qfigdef{p2_phy_s3_q06_1}{TOPS - Test 02 (28 Aug 2022)-img/p2_phy_s3_q06_1.png}{82.0}{0}
\raggedcolumns
\makeatletter\@beginparpenalty=10000\makeatother
\newcolumntype{L}[1]{>{\raggedright\arraybackslash}p{#1}}
\begin{document}
\begin{center}{\Large\bfseries TOPS --- Test 02 (28 Aug 2022)}\end{center}
\begin{multicols}{2}
\raggedcolumns
\papersec{1}
\partsec{1}{PHYSICS}
\sect{SECTION-I(i)}
\qn{1} A particle is moving such that its position vector changes with time as
\par\noindent\hfill$\begin{aligned}\vec r&=\left(36t^2-4t^3-96t\right)\hat i+\left(12t^3-108t^2+288t-17\right)\hat j\\&\quad+\left(51+3t^3-27t^2+72t\right)\hat k\end{aligned}$\hfill\null\par
What is the distance travelled by the particle in the first six seconds?
\begin{qoptions}
\opt{A}{216}
\opt{B}{384}
\opt{C}{468}
\opt{D}{572}
\end{qoptions}
\qn{2} In which of the following cases is the work done non-zero?
\begin{qoptions}
\opt{A}{Work done by normal of ladder on you when you climb a ladder}
\opt{B}{Work done by friction on your shoes when you run in a race}
\opt{C}{Work done by a pencil on a paper when you write on the paper}
\opt{D}{None of the above}
\end{qoptions}
\qn{3} A rope of length $L$ and $M$ is being pulled on a rough floor by a constant horizontal force $F=Mg$. The force is acting at one end of the rope in the same direction as the length of the rope. The coefficient of kinetic friction between rope and floor is $1/2$. Then, tension at the mid-point of the rope is
\begin{qoptions}
\opt{A}{$\dfrac{Mg}4$}
\opt{B}{$\dfrac{2Mg}5$}
\opt{C}{$\dfrac{Mg}8$}
\opt{D}{$\dfrac{Mg}2$}
\end{qoptions}
\qn{4} All pulleys and threads are ideal. $C$ is a point on the thread as shown. When the system is released from rest, at that instant
\qfig{p1_phy_s1_q04_1}
\begin{qoptions}
\opt{A}{Acceleration of point $C$ is $5g/2$ upwards}
\opt{B}{Acceleration of block $B$ is $2g$ upwards}
\opt{C}{Acceleration of block $A$ is $g/2$ downwards}
\opt{D}{Acceleration of point $C$ is $g$ downwards}
\end{qoptions}
\qn{5} Two identical solid cylinders of radius $r$ are placed touching along their lengths on a horizontal surface. A third smooth cylinder of same material but twice the radius of that of the cylinders is placed lengthwise on them so that the system remains at rest. If all three cylinders have the same length, then minimum value of the coefficient of friction between smaller cylinders and the surface is :
\qfig{p1_phy_s1_q05_1}
\begin{qoptions}
\opt{A}{$\dfrac1{\sqrt2}$}
\opt{B}{$\dfrac13$}
\opt{C}{$\dfrac1{3\sqrt2}$}
\opt{D}{None of these}
\end{qoptions}
\qn{6} Consider a straight line passing through two points $A$ and $B$. Position vectors of point $A$ and $B$ are given by $\vec r_A$ and $\vec r_B$ respectively. There is another point $C$ with position vector $\vec r_C$. If a perpendicular is dropped on the line passing through $A$ and $B$ from the point $C$. Then foot of perpendicular on straight line is given by point $D$ with a position vector $\vec r_D=x_0\hat i+y_0\hat j+z_0\hat k$ if $\vec r_A=4\hat i+7\hat j+\hat k$, $\vec r_B=3\hat i+5\hat j+3\hat k$ and $\vec r_C=\hat i+3\hat k$. Then the value of $[x_0+z_0-y_0]$ will be?
\begin{qoptions}
\opt{A}{12}
\opt{B}{5}
\opt{C}{8}
\opt{D}{0}
\end{qoptions}
\sect{SECTION-I(ii)}
\qn{7} All the pulleys are frictionless and massless. String is inextensible and massless. A man of Mass $2M$ is climbing up with acceleration of $a_0$ relative to the point he is holding on the string. A block of mass $M$ is connected with the help of string and pulleys as shown. If tension in the string which man is holding, is $T$. Acceleration of mass $M$ is $a_1$ upwards and acceleration of man with respect to ground is $A_0$ upwards. Then choose the correct options.
\qfig{p1_phy_s2_q07_1}
\begin{qoptions}
\opt{A}{$a_1=\dfrac{3g+4a_0}9$}
\opt{B}{$A_0=\dfrac{a_0-6g}9$}
\opt{C}{$T=\dfrac{Ma_0+3Mg}9$}
\opt{D}{$T=\dfrac{2Ma_0+6Mg}9$}
\end{qoptions}
\qn{8} An iron sphere weighing 10N rests in a V shaped smooth trough whose sides form an angle of $60^\circ$ as shown in the figure. Then the reaction forces are ($g=10$ m/s$^2$)
\qfig{p1_phy_s2_q08_1}
\begin{qoptions}
\opt{A}{$R_A=10$N \& $R_B=0$ in case (i)}
\opt{B}{$R_A=10$N \& $R_B=10$N in case (ii)}
\opt{C}{$R_A=\dfrac{20}{\sqrt3}$N \& $R_B=\dfrac{10}{\sqrt3}$N in case (iii)}
\opt{D}{$R_A=10$N \& $R_B=10$N in all the three cases}
\end{qoptions}
\qn{9} Which of the following statements is correct?
\begin{qoptions}
\opt{A}{Uniform circular motion may be a uniform motion or a uniformly accelerated motion}
\opt{B}{Non-linear motion cannot be a uniformly accelerated motion as there must be some centripetal acceleration which keeps on changing in direction}
\opt{C}{For an object in circular motion with only two forces acting on it, if one of the forces is constant in magnitude and direction, the other force must be varying in magnitude as well as in direction}
\opt{D}{None of the above}
\end{qoptions}
\qn{10} A small toy train is moving with a constant speed of 10 m/s in a circle of radius $\dfrac{16}{\pi}$ m. The plane of the circle lies in horizontal $x$--$y$ plane. At time $t=0$, train is at point $P$ and moving in counter--clockwise direction. At this instant, a stone is thrown from the train with speed 10 m/s relative to train towards negative $x$--axis at an angle of $37^\circ$ with vertical $z$--axis.

(Take $g=10$ m/s$^2$, $\sin37^\circ=\dfrac35$) (Given : centre of circle at origin $0,0$)
\qfig{p1_phy_s2_q10_1}
\begin{qoptions}
\opt{A}{Velocity of particle at highest point will become $\left(-6\hat i+10\hat j\right)$ m/s.}
\opt{B}{Velocity of particle at highest point will become $\left(-8\hat i+10\hat j\right)$ m/s.}
\opt{C}{Coordinates at highest point are, (0.3m, 7.0m, 3.2 m).}
\opt{D}{Coordinates at highest point are, (0.3m, 8.0m, 3.2m).}
\end{qoptions}
\qn{11} A sledge moving over a smooth horizontal surface of ice at a velocity $v_0$ drives out on a horizontal road and comes to a halt as shown. The sledge has a length $\ell$, mass $m$ and friction between sledge and road is $\mu$.
\qfig{p1_phy_s2_q11_1}
\begin{qoptions}
\opt{A}{No work is done by the friction to switch the sledge from ice to the road.}
\opt{B}{A work of $\dfrac12\mu mg\ell$ is done against friction while sledge switches completely on to road.}
\opt{C}{The distance covered by the sledge on the road is $\left(\dfrac{v_0^2}{2\mu g}-\dfrac\ell2\right)$.}
\opt{D}{Total distance moved by the sledge before stopping is $\left(\dfrac{v_0^2}{2\mu g}+\dfrac\ell2\right)$.}
\end{qoptions}
\qn{12} Three large box are kept stationary over one other as shown in figure. A horizontal force 20 N acts on middle block as shown in figure. Then : ($g=10$ m/s$^2$)
\qfig{p1_phy_s2_q12_1}
\begin{qoptions}
\opt{A}{Acceleration of lowest block is 4.5 m/s$^2$.}
\opt{B}{Acceleration of upper most block is 2 m/s$^2$.}
\opt{C}{Friction between middle and lower block is 12 N.}
\opt{D}{Friction between middle and upper block is 2 N.}
\end{qoptions}
\sect{SECTION-II}
\qn{1} Point `$A$' moves uniformly with speed $v_1=20$ m/s, so that vector $\vec v_1$ is continuously `aimed' at point `$B$' which in turn moves rectilinearly and uniformly with vector $v_2=10$ m/s along the path $P\to Q$ as shown in the figure. Find the time (in sec) when the points $A$ and $B$ converge.
\qfig{p1_phy_s3_q01_1}
\qn{2} A particle is projected up an incline at $t=0$ (inclination $=53^\circ$) with speed 14 m/s and at an angle $30^\circ$ from the incline. At the same moment a box starts sliding down, the box has 2 small hole at height of 2 m on opposite walls. The particle enters form 1$^{\text{st}}$ hole and leave through the other. Find the width of the box.
\qfig{p1_phy_s3_q02_1}
\qn{3} A uniform chain of mass per unit length $\lambda$ is sliding down the hill with constant speed $u$. Dimension of the hill are $h$ and $d$ as shown. Find the momentum (in kg -- m/s) of the chain at the instant shown. (Given $\lambda=1$ kg/m, $d=3$m, $h=4$m, and $u=1$m/s)
\qfig{p1_phy_s3_q03_1}
\qn{4} In the arrangement shown, pulleys and string are massless and all surfaces are frictionless. The block and the wedge have same mass $M$. The string is being pulled with a constant force $F$. Relative acceleration between the block and the wedge has a magnitude of $\dfrac{KF}M$. Find $K$.
\qfig{p1_phy_s3_q04_1}
\qn{5} The spring block system lies on a smooth horizontal surface. The free end of the spring is being pulled towards right with constant speed $v_0=2$ m/s. At $t=0$ sec, the spring of constant $k=100$ N/cm is unstretched and the block has a speed 1 m/s to left. Find the maximum extension of the spring in cm.
\qfig{p1_phy_s3_q05_1}
\qn{6} A toy car is moving on a circular path of radius 10 m on an inclined plane of inclination $30^\circ$. The coefficient of friction between the tyres and the inclined plane is $\mu=\dfrac2{\sqrt3}$. The maximum constant speed of car, such that car does not slip on the inclined plane is $v_0$ m/s. Assume that width of the car is negligible with respect to radius of circular path. Find the value of $v_0$. ($g=10$ m/s$^2$)
\qfig{p1_phy_s3_q06_1}
\columnbreak
\partsec{2}{CHEMISTRY}
\sect{SECTION-I(i)}
\qn{1} How many hydrogen atoms are present in 0.1 mole of mohr's salt ($\mathrm{FeSO_4.(NH_4)_2SO_4.6H_2O}$) ?
\begin{qoptions}
\opt{A}{$10\times N_A$}
\opt{B}{$14\times N_A$}
\opt{C}{$2\times N_A$}
\opt{D}{$0.8\,N_A$}
\end{qoptions}
\qn{2} A sample of water contains 10 gm hydrogen atom. The weight of oxygen present in the given sample is
\begin{qoptions}
\opt{A}{160 gm}
\opt{B}{320 gm}
\opt{C}{240 gm}
\opt{D}{80 gm}
\end{qoptions}
\qn{3} How many mole of protons are present in 90g of H$_2$O ?
\begin{qoptions}
\opt{A}{50 mole}
\opt{B}{60 mole}
\opt{C}{70 mole}
\opt{D}{40 mole}
\end{qoptions}
\qn{4} Total number of faces and edges in IF$_7$ polyhedron respectively are :
\begin{qoptions}
\opt{A}{10, 14}
\opt{B}{5, 15}
\opt{C}{15, 10}
\opt{D}{10, 15}
\end{qoptions}
\qn{5} Which of the following is the correct order of angle between two oxygen atoms linked to central atom in the given molecules

I. CO$_2$\quad II. XeO$_4$\quad III. XeO$_2$F$_4$\quad IV. SO$_3$
\begin{qoptions}
\opt{A}{I $>$ II $>$ III $>$ IV}
\opt{B}{I $>$ III $>$ II $>$ IV}
\opt{C}{III $>$ I $>$ IV $>$ II}
\opt{D}{III $=$ I $>$ IV $>$ II}
\end{qoptions}
\qn{6} Select the correct match(es) :-
\begin{qoptions}
\opt{A}{Monoclinic sulphur : Extended covalently bonded solid}
\opt{B}{Black phosphorous : discrete molecular structure}
\opt{C}{Dry ice : ionic solid}
\opt{D}{C$_{60}$ : discrete molecular structure}
\end{qoptions}
\sect{SECTION-I(ii)}
\qn{7} 124 u of P$_4$ will contain
\begin{qoptions}
\opt{A}{4 $N_A$ atoms of phosphorus}
\opt{B}{4 atoms of phosphorus}
\opt{C}{1 molecule of phosphorus}
\opt{D}{$N_A$ molecules of phosphorus}
\end{qoptions}
\qn{8} Which is/are correct statements ?
\begin{qoptions}
\opt{A}{Molar volume of H$_2$O($\ell$) is almost equal to 18 ml at 1 atm pressure and 298 K temperature}
\opt{B}{Molar volume of N$_2$ gas is equal to 22.4 L at 1 atm pressure and 298 K temperature}
\opt{C}{Molar volume of O$_2$ gas is equal to 22.7 L at STP}
\opt{D}{Molar volume of NH$_3$ gas is equal to 22.4 L at 1 atm and 273 K.}
\end{qoptions}
\qn{9} Which is/are correct statements about 10 mol of Na$_2$SO$_4$.7H$_2$O.
\begin{qoptions}
\opt{A}{$6.022\times10^{24}$ atoms of sulphur present}
\opt{B}{$7\times6.022\times10^{24}$ molecules of H$_2$O present}
\opt{C}{176 amu of oxygen present}
\opt{D}{46 gm of sodium present}
\end{qoptions}
\qn{10} Consider the following molecules.

(P) P$_4$O$_8$

(Q) C$_3$O$_2$

(R) N$_3$H

(S) (CN)$_2$

Select the incorrect statement(s).
\begin{qoptions}
\opt{A}{Hybridisation of atleast 2 atoms in two molecules is sp$^3$}
\opt{B}{Hybridisation of atleast 2 atoms in two molecules is sp}
\opt{C}{One molecule contain 3p$\pi$-3d$\pi$ bond}
\opt{D}{Only two molecules contain one 2p$\pi$-2p$\pi$ bond}
\end{qoptions}
\qn{11} Which of the following is/are correctly matched for the given species?

BeCl$_2$, I(CN)$_2^-$, IF$_7$, CH$_2$Cl$_2$, SiF$_4$, ICl$_2^-$, H$_2$O, SF$_4$, BF$_3$, XeF$_3^+$, I$_3^+$
\begin{qoptions}
\opt{A}{Polar species $=5$}
\opt{B}{V-shape species $=1$}
\opt{C}{Planar species $=6$}
\opt{D}{Number of species in which atleast one bond angle is $180^\circ=4$}
\end{qoptions}
\qn{12} To attain the maximum stabilisation through H bonding, the different types of structures such as chain structure, sheet structure and three dimentional structure are formed.

Select the correct combination.
\begin{qoptions}
\opt{A}{Ice (3d - structure)}
\opt{B}{HCN (Zig-Zag)}
\opt{C}{Boric acid (Sheet structure)}
\opt{D}{CH$_3$COOH (Dimeric)}
\end{qoptions}
\sect{SECTION-II}
\qn{1} For a gaseous substance a graph between various mole was plotted against the respective mass(g) of the substance as given below. What is the volume ($\ell$) of 5.6g of the substance at 1atm and 273K.
\qfig{p1_chem_s3_q01_1}
\qn{2} Adipic acid HOOC(CH$_2$)$_4$COOH is used in making J nylon. Calculate the weight of hydrogen atom in gram present in 58.4 g of adipic acid.
\qn{3} If iron (Fe 56) is the only metal in a protein and rest are non-metals, how many iron atoms are present in one molecule of protein if the molar mass of protein is 64000 gm/mol and it contains 0.35\% iron by mass? (Atomic mass of Fe $=56$)
\qn{4} Find the number of compounds in which covalently bonded cyclic structure observed .

Na$_2$S$_4$O$_6$, Black phosphorous, Graphite, Rhombic sulphur, H$_6$P$_4$O$_{13}$, P$_4$S$_{10}$, S$_3$O$_9$, S$_2$F$_{10}$
\qn{5} The number of molecules or ions in which tetrahedral units connected by oxy-linkage

S$_3$O$_9$, B$_3$O$_6^{3-}$, S$_4$O$_6^{2-}$, Na$_2$B$_4$O$_7$.10H$_2$O, H$_2$S$_2$O$_6$, N$_2$O$_3$(sym), P$_4$O$_{10}$, H$_2$S$_2$O$_8$, H$_2$S$_3$O$_6$
\qn{6} Find the multiplication of $x$ and $y$.

if $x=$ sum of $90^\circ$ bond angles present in $a,b,c$ species

$y=$ sum of maximum number of bonds having equal bond length in $a,b,c$ species

where,

$a=$ Cationic part of XeF$_6$(s)

$b=$ Anionic part of I$_2$Cl$_6$($\ell$)

$c=$ Anionic part of PCl$_5$(s)
\partsec{3}{MATHEMATICS}
\sect{SECTION-I(i)}
\qn{1} Sum of 10 terms of series 2, 6, 14, 30, 62, \ldots, is -
\begin{qoptions}
\opt{A}{$2^{24}$}
\opt{B}{$2^{12}-2$}
\opt{C}{$2^{24}-12$}
\opt{D}{$2^{12}-24$}
\end{qoptions}
\qn{2} The complete interval of values of $x$ in $\left(-\dfrac\pi2,\dfrac\pi2\right)$ satisfying the inequations $\cos x.\cos2x>-\dfrac14$ and $3\cos x>1$, is
\begin{qoptions}
\opt{A}{$\left(0,\dfrac\pi3\right)$}
\opt{B}{$\left(-\dfrac\pi2,\dfrac\pi4\right)$}
\opt{C}{$\left(-\dfrac\pi3,\dfrac\pi3\right)$}
\opt{D}{$\left(-\dfrac\pi3,\dfrac{5\pi}{12}\right)$}
\end{qoptions}
\qn{3} If $\log_{10}25=a$, $\log_{10}225=b$ then the value of $\log_{10}\left(\dfrac19\right)^2+\log_{10}\left(\dfrac1{2250}\right)$ is
\begin{qoptions}
\opt{A}{$2a-3b-1$}
\opt{B}{$2a+3b$}
\opt{C}{$2a+3b+1$}
\opt{D}{$2a+3b-1$}
\end{qoptions}
\qn{4} Let $p(x)=ax^2+bx+c$ such that $a>0$, $a+b+c\ge0$, $a-b+c\ge0$, $a-c\ge0$ and $b^2>4ac$, then the real roots $x_1,x_2$ of $p(x)=0$
\begin{qoptions}
\opt{A}{must lie in $[-1,1]$}
\opt{B}{must lie in $[1,\infty)$}
\opt{C}{must lie in $(-\infty,-1]$}
\opt{D}{must lie in $[0,1]$}
\end{qoptions}
\qn{5} The coefficient of $x^{17}$ in $\displaystyle\prod_{i=1}^{18}(x-i)$ is
\begin{qoptions}
\opt{A}{$\dfrac{171}2$}
\opt{B}{342}
\opt{C}{$-171$}
\opt{D}{684}
\end{qoptions}
\qn{6} If $P=\left(\tan3^{n+1}\theta-\tan\theta\right)$ and $Q=\displaystyle\sum_{r=0}^{n}\dfrac{\sin3^r\theta}{\cos3^{r+1}\theta}$ then
\begin{qoptions}
\opt{A}{$P=3Q$}
\opt{B}{$P=2Q$}
\opt{C}{$2P=Q$}
\opt{D}{$3P=Q$}
\end{qoptions}
\sect{SECTION-I(ii)}
\qn{7} Which of the following is/are true ?
\begin{qoptions}
\opt{A}{$(\tan x)^{\ell n(\sin x)}>(\cot x)^{\ell n(\sin x)}\ \forall x\in\left(0,\dfrac\pi4\right)$}
\opt{B}{$4^{\ell n(\operatorname{cosec}x)}>5^{\ell n(\operatorname{cosec}x)}\ \forall x\in\left(0,\dfrac\pi2\right)$}
\opt{C}{$\left(\dfrac12\right)^{\ell n(\cos x)}<\left(\dfrac13\right)^{\ell n(\cos x)}$, $\forall x\in\left(0,\dfrac\pi2\right)$}
\opt{D}{$2^{\ell n(\tan x)}>2^{\ell n(\sin x)}$, $\forall x\in\left(0,\dfrac\pi2\right)$}
\end{qoptions}
\qn{8} If $n$ is a natural number, then the value of $\displaystyle\sum_{r=1}^{n}r^3$ can be -
\begin{qoptions}
\opt{A}{225}
\opt{B}{441}
\opt{C}{36}
\opt{D}{25}
\end{qoptions}
\qn{9} Let $x,y,z$ are real numbers such that $x,y,z\in\left(0,\dfrac\pi2\right)$ and $x>y>z$. If $x+y+z=\dfrac{3\pi}4$, $\tan x+\tan y+\tan z=5$ and $\tan x\tan y\tan z=1$, then
\begin{qoptions}
\opt{A}{$x=\dfrac{5\pi}{12}$}
\opt{B}{$y=\dfrac\pi4$}
\opt{C}{$y=\dfrac\pi{12}$}
\opt{D}{$z=\dfrac\pi{12}$}
\end{qoptions}
\qn{10} Let $A$, $G$, $H$ are arithmetic, geometric and harmonic mean between two given positive numbers $x$ and $y$. Then -
\begin{qoptions}
\opt{A}{$A^2$, $AG$, $AH$ are in G.P.}
\opt{B}{Quadratic $Ax^2+Gx+H=0$ has imaginary roots}
\opt{C}{Quadratic $Ax^2+Gx+H=0$ has real roots}
\opt{D}{$A\ge G\ge H$}
\end{qoptions}
\qn{11} If $S$ is solution set of all values of $x$ for which $2\sin t=\dfrac{3x+1}{x-1}$, then
\begin{qoptions}
\opt{A}{number of integral values in $S$ is 4}
\opt{B}{number of integral values in $S$ is 5}
\opt{C}{Sum of all integral values in $S$ is $-6$}
\opt{D}{Sum of all integral values in $S$ is $-5$}
\end{qoptions}
\qn{12} If $S$ is sum of solutions \& $P$ is product of solutions of the equation $2^{\log_2(x^2-x)}-4x+5^{\log_{10}6}.2^{\log_{10}6}=0$, then -
\begin{qoptions}
\opt{A}{$S+P=11$}
\opt{B}{$S^2=16$}
\opt{C}{$P^2=25$}
\opt{D}{$S^2=25$}
\end{qoptions}
\sect{SECTION-II}
\qn{1} Let $a,b,c$ be distinct positive numbers in A.P. with minimum common difference $d$, where $d\in I$ such that their G.M. is $b^{2/3}$. If $c>a+2$, then the value of $(2b-1)^2$ is
\qn{2} The value of $\dfrac{\operatorname{cosec}10^\circ+\sec20^\circ}{\operatorname{cosec}20^\circ+\sec10^\circ}$ is equal to
\qn{3} If $x^2-5x+4=0$ and $x^2-px+12=0$ has exactly one common root and sum of possible values of $P$ is $\lambda$, then $\lambda$ is
\qn{4} Number of solution(s) of the equation $27^{\sin^2x}+27^{\cos^2x}=12$ in the interval $[-\pi,\pi]$ is
\qn{5} If the solution set of inequality $\left(\cos x+\dfrac{\sqrt3}2\right)\left(\cos x-\dfrac12\right)\le0$ in $[0,2\pi]$ is $\left[\dfrac{\alpha\pi}6,\dfrac{\beta\pi}6\right]\cup\left[\dfrac{\gamma\pi}6,\dfrac{\delta\pi}6\right]$, $\alpha,\beta,\gamma,\delta\in I^+$, then the value of $|\beta-\alpha+\delta-\gamma|$ is
\qn{6} If $\log_7^2(x+7)+\log_7^2(3+\log_2y)=0$, then value of $|x|+4y$ is
\papersec{2}
\partsec{1}{PHYSICS}
\sect{SECTION-I}
\qn{1} An air bubble under explosion inside water oscillates with a time period given by $T=P^aE^bd^c$ where `$P$' is pressure of air, `$E$' is energy of explosion, `$d$' density of liquid then :-
\begin{qoptions}
\opt{A}{$a=-5/6$}
\opt{B}{$b=-0.4\,a$}
\opt{C}{$c=0.5$}
\opt{D}{None of these}
\end{qoptions}
\qn{2} Bead $A$ can slide on a smooth horizontal rod and block $B$ move vertically downward. String connected between them and Pulley are ideal. Block is released gently and at a certain instant (not at the time of releasing) the velocity of ring is $v_1$ and block is $v_2$, At that instant ring is at horizontal distance `$x$' from pulley and vertical distance is `$y$' then mark the correct option(s):
\qfig{p2_phy_s1_q02_1}
\begin{qoptions}
\opt{A}{$v_2=v_1\cos\theta$}
\opt{B}{$v_1=v_2\cos\theta$}
\opt{C}{$a_2=a_1\cos\theta$, where $a_2$ and $a_1$ represent the acceleration of block and ring respectively at that instant.}
\opt{D}{Angular velocity of the ring with respect to an observer on the pulley is $\dfrac{v_1\sin\theta}{\sqrt{x^2+y^2}}$.}
\end{qoptions}
\qn{3} A projectile of mass $m$ is fired at an angle $\theta_0$ with horizontal and initial speed $V_0$, as shown. If the air develops a drag force (resistance force) on the projectile which is proportional to its instantaneous velocity i.e. $F=-KV$ where $K$ is a positive constant of appropriate dimensions. (Assume that gravitational acceleration is constant throughout the motion i.e. $g$ m/s$^2$)
\qfig{p2_phy_s1_q03_1}
\begin{qoptions}
\opt{A}{The time of ascent is greater than the time of decent.}
\opt{B}{Vertical ($Y$-component) component of velocity of projectile as a function of time is $\dfrac mK\left[\left(\dfrac Km V_0\sin\theta_0-g\right)e^{\frac{-K}mt}+g\right]$}
\opt{C}{Horizontal ($X$-component) component of velocity of projectile as a function of time is $V_0\cos\theta_0e^{-\frac kmt}$}
\opt{D}{Vertical ($Y$-component) component of velocity of projectile as a function of time is $\dfrac mK\left[\left(\dfrac Km V_0\sin\theta_0+g\right)e^{\frac{-2K}mt}+g\right]$}
\end{qoptions}
\qn{4} A wedge of mass $m$ is kept on a smooth table and its inclined surface is also smooth. A small block of mass $m$ is projected from the bottom along the incline surface with velocity $u$. Assume that the block remains on the incline and take $\theta=45^\circ$, $g=10$ m/s$^2$.
\qfig{p2_phy_s1_q04_1}
\begin{qoptions}
\opt{A}{Magnitude of acceleration of block is more than magnitude of acceleration of wedge.}
\opt{B}{The block strikes the table normally.}
\opt{C}{Radius of curvature of the path of the block when it is at the highest point is $u^2/16g$.}
\opt{D}{Radius of curvature of the path of the block when it is at the highest point is $3u^2/16g$.}
\end{qoptions}
\qn{5} In the given situation, coefficient of friction between the blocks $A$ and $B$ is $\mu$ and there is no friction between ground and block $B$, acceleration of $A$ is $a_A$ and acceleration of $B$ is $a_B$. A force is acting on either of the blocks, then mark the correct option(s) : (Take $g=10$ m/s$^2$)
\qfig{p2_phy_s1_q05_1}
\begin{qoptions}
\opt{A}{If acceleration of $A$ and $B$ is in same direction and its magnitude is 2 m/s$^2$ and 3 m/s$^2$ respectively then external force $F$ must act on block $B$.}
\opt{B}{If acceleration of $A$ and $B$ is in same direction and its magnitude is 2 m/s$^2$ and 3 m/s$^2$ respectively, mass of $A$ is 20 kg and mass of $B$ is 50 kg, then friction coefficient between the bodies is 0.4}
\opt{C}{If mass of $A$ is 10 kg and mass of $B$ is 25 kg, $\mu=0.4$ and $F=30$ N is applied on block $A$ then friction between the blocks is static in nature.}
\opt{D}{If mass of $A$ is 10 kg and mass of $B$ is 25 kg, $\mu=0.4$ and $F=250$ N is applied on block $A$ then $a_A=(50/7)$ m/s$^2$.}
\end{qoptions}
\qn{6} When a particle is in motion along a curved path, its velocity at a point $A$ is along the tangent to the path at point $A$ as shown in the first figure. We may split this velocity into two components; one along position vector $OA$ and call it radial velocity and another along the perpendicular to $OA$ and call it transverse velocity. Same can be done to acceleration at $A$ and call the components radial acceleration and transverse acceleration.
\qfig{p2_phy_s1_q06_1}
A particle is executing uniform circular motion of angular velocity $\omega$ in $x$--$y$ plane with radius $R$ center $(2R,0)$. When it is at point $A$, answer the correct option(s).
\qfig{p2_phy_s1_q06_2}
\begin{qoptions}
\opt{A}{Ratio of angular velocity to the radial velocity of the particle about origin is $1/R$.}
\opt{B}{Ratio of angular velocity to the radial velocity of the particle about origin is $\dfrac1{2\sqrt5R}$.}
\opt{C}{Net acceleration of particle about origin is $R\omega^2$.}
\opt{D}{To find angular acceleration about the origin in the given question we can apply $a_t=\alpha r$, where $a_t$ is transverse acceleration about point $O$ and $r$ is distance between point $O$ and $A$.}
\end{qoptions}
\sect{SECTION-II}
\qn{1} A ball is projected vertically upwards with a velocity 20 m/s returns back on the ground with a velocity 16 m/s. if the air resistance is proportional to the speed of ball find the air time of the ball. Acceleration of free ball is 10 SI unit.
\qn{2} An elastic cord of force constant $K=0.25$ N/m$^2$ and natural length 10 cm is hanging from a ceiling. An insect of mass 5 g is sitting at the lower end of the cord and the system is in equilibrium. The insect start climbing up the spring toward the ceiling. Assume that there is no slipping and insect is moving slowly. Length of the cord, when insect is $1/3^{\text{rd}}$ of its original distance from ceiling, is $L$ cm, then $L$ is equal to (given that $g=10$ m/s$^2$)
\qfig{p2_phy_s2_q02_1}
\qn{3} Initially system is in equilibrium, Now mass $m$ is given velocity $v$ in horizontal direction as shown. Find minimum velocity $v$ in m/s Such that either of block $A$ or $B$ loose contact with ground. given $m_A=m_B=5$kg : ($m=1$ kg, $g=10$ m/s$^2$)
\qfig{p2_phy_s2_q03_1}
\qn{4} Find the magnitude of acceleration (in m/s$^2$) of pulley of mass 1 kg as shown in figure. Other pulley is massless and string is ideal. (Take : $g=9.8$ m/s$^2$)
\qfig{p2_phy_s2_q04_1}
\qn{5} In the figure all pulleys are massless and frictionless and strings are light and inextensible. The system consists of infinite number of pulleys and blocks (Blocks are identical and having mass 2 kg). Pulley $P_1$ is moved to left with a constant acceleration of 3 m/s$^2$, find the acceleration of block 1. Assume the strings to remain horizontal.
\qfig{p2_phy_s2_q05_1}
\qn{6} Attraction between two uncharged conducting parallel plates placed in vacuum is described by Casimir effect. This force per unit area of each plate depends on plank's constant ($h$), speed of light ($c$) and separation between the plates ($r$). Using dimensional analysis, the dimensionless proportionality constant is `$k$'. If the force acting on $1\times1$ cm plates separated by 1 $\mu$m is 0.013 dyne, calculate the value of constant $k$ is $P\times10^{-3}$. The value of $P$ is
\sect{SECTION-III}
\qn{1} A particle rotates in a circle with angular speed 18 rad/sec. A retarding force decelerates it such that angular deceleration is always proportional to square root of angular velocity. Find the magnitude of mean angular velocity (in rad/sec) of the particle averaged over the whole time of rotation.
\qn{2} A racing car driver drives his car on a flat circular track of radius $25/3$ m and a coefficient of friction of 0.5. He drives the car in such a manner that he may attain the maximum possible velocity on the track in a minimum possible time. The magnitude of his tangential acceleration (in SI unit) at an instant when his speed is 5 m/s is
\qn{3} A rubber ball is released from a height of 5 m above the floor. It bounces back repeatedly, always rising to 0.81 times of the height through which it falls. The average speed of the ball is $v$ (in m/s), then value of $2v$ to nearest integer is : (Take $g=10$ ms$^{-2}$)
\qn{4} Two blocks of equal mass, 4 kg each, are connected to two ends of a massless string passing over a massless pulley. On one side of the string there is a bead of mass 2kg as shown. When the system is released from rest the bead continues to remain at rest while the two blocks accelerate. If the acceleration of the blocks in SI unit is $p$ then value of $2p$ is. (Take $g=10$m/s$^2$)
\qfig{p2_phy_s3_q04_1}
\qn{5} A fireman fires his shot from origin aiming to monkey ($M$), location at $(12, 16)$, who falls at the time of shot. At the instant shot has passes its highest position, it is located on line $x=8$ m, the monkey is located at position $(12,n/3)$. Find the value of $n/4$. (Coordinates are in m)
\qn{6} There is a U shaped container of Uniform cross sectional area $1$ cm$^2$, It is suspended by two ideal strings $S_1$ and $S_2$ and ideal Spring of force constant $K=10$ N/m. An unknown liquid of density is poured in it, assume that density of liquid remain same throughout the volume of the container. It is observed that when same liquid is poured into the container, the level of liquid remain unchanged with respect to surrounding person at ground frame. Given that $g=10$ m/s$^2$ and $\theta=30$ degree. Density of liquid is (in SI unit) $d\times10^3$, then $d$ is :
\qfig{p2_phy_s3_q06_1}
\partsec{2}{CHEMISTRY}
\sect{SECTION-I}
\qn{1} Select the correct statement(s) for $\mathrm{(NH_4)_3PO_4}$.
\begin{qoptions}
\opt{A}{Ratio of number of oxygen atoms to number of hydrogen atoms is 1 : 3}
\opt{B}{Ratio of number of cations to number of anions is 3 : 1}
\opt{C}{Ratio of number of gm-atoms of nitrogen to gm-atoms of oxygen is 3 : 2}
\opt{D}{Total number of atoms in one mole of $\mathrm{(NH_4)_3PO_4}$ is 20}
\end{qoptions}
\qn{2} Which of the following has a mass of 20 g. [Use $N_A=6\times10^{23}$]

[Atomic mass : Ca $=40$, S $=32$, O $=16$, He $=4$, C $=12$]
\begin{qoptions}
\opt{A}{$3\times10^{23}$ Ca$^{2+}$ ion}
\opt{B}{5.6 L SO$_3$ gas at 1 atm and 273 K}
\opt{C}{Mixture of 1 mol He gas and 1 mol CH$_4$ gas}
\opt{D}{$1.2\times10^{23}$ molecule of CaCO$_3$.}
\end{qoptions}
\qn{3} Identify the correct statements among the following
\begin{qoptions}
\opt{A}{0.5 gm-atom of oxygen contains $\dfrac{N_A}4$ oxygen molecules}
\opt{B}{Total electrons present in 0.5 g-atom of oxygen is $4N_A$}
\opt{C}{Mass of $3.01\times10^{23}$ molecules of ammonia is 8.5 g}
\opt{D}{Number of atoms present in 11.2 litres of chlorine at 1 atm \& 273K is $6.022\times10^{23}$}
\end{qoptions}
\qn{4} Consider the following statements :

$S_1=$ molecule/ion is non-planar

$S_2=$ all bond lengths are same in molecule/ion

$S_3=$ non polar bond(s) is/are present in molecule/ion

Select correct option(s) considering above statement(s).
\begin{qoptions}
\opt{A}{Only $S_1$ and $S_3$ are correct for fullerene (C$_{60}$)}
\opt{B}{Only $S_2$ and $S_3$ is correct for Azide ion}
\opt{C}{Only $S_1$ is correct for hydrogen peroxide}
\opt{D}{Only $S_2$ and $S_3$ are incorrect for tetrathionate ion}
\end{qoptions}
\qn{5} Select the correct option?
\begin{qoptions}
\opt{A}{$d_{\mathrm{P-O}}\,(\mathrm{POF_3})>d_{\mathrm{P-O}}\,(\mathrm{PO_3F^{2-}})$}
\opt{B}{$d_{\mathrm{P-O}}\,(\mathrm{PO_4^{3-}})>d_{\mathrm{P-O}}\,(\mathrm{PO_3F^{2-}})$}
\opt{C}{$d_{\mathrm{P-F}}\,(\mathrm{POF_3})>d_{\mathrm{P-F}}\,(\mathrm{PO_3F^{2-}})$}
\opt{D}{$d_{\mathrm{P-O}}\,(\mathrm{POCl_3})>d_{\mathrm{P-O}}\,(\mathrm{POF_3})$}
\end{qoptions}
\qn{6} Process-(1)

PCl$_3+$ BBr$_3\xrightarrow[\text{Product(1)}]{\text{adduct of PCl}_3\text{ and BBr}_3}$

Process-(2)

PCl$_3+$ N(CH$_3$)$_3\xrightarrow[\text{Product(2)}]{\text{adduct of PCl}_3\text{ and N(CH}_3)_3}$

Select INCORRECT statement(s) considering process-1 and 2.
\begin{qoptions}
\opt{A}{In both processes no atom undergoes changes in hybridisation}
\opt{B}{No atom is forming bond in excited state in given products}
\opt{C}{In both process phosphorus atom donate lone pair}
\opt{D}{In product-(2) total 4 tetrahedral units present}
\end{qoptions}
\sect{SECTION-II}
\qn{1} Let $x$(g) be mass of CO$_2$ that must be taken so that it contains same number of O atoms as there number of H atoms in 9 g H$_2$O. Find $x$.
\qn{2} $2.4\times10^y$ molecules of water are present in 252 mg of H$_2$C$_2$O$_4$.2H$_2$O then value of $y$ is-

[Given $N_A=6\times10^{23}$]
\qn{3} A compound possesses 8\% sulphur by mass. The least molecular mass of the compound is
\qn{4} How many orders are correct :

(i) H--F $<$ H--Cl $<$ H--Br $<$ H--I (Bond length)

(ii) H--F $<$ H--Cl $<$ H--Br $<$ H--I (Acidic strength)

(iii) H--I $<$ H--Br $<$ H--Cl $<$ H--F (Bond strength)

(iv) H--F $>$ H--Cl $>$ H--Br $>$ H--I (Thermal stability)

(v) H--F $>$ H--I $>$ H--Br $>$ H--Cl (Boiling point)

(vi) F$_2<$ Cl$_2<$ Br$_2<$ I$_2$ (Bond length)

(vii) F$_2<$ Cl$_2<$ Br$_2<$ I$_2$ (Bond energy)

(viii) F$_2<$ Cl$_2<$ Br$_2<$ I$_2$ (Boiling point)

(ix) F $>$ Cl $>$ Br $>$ I (Electron affinity)

(x) F $>$ Cl $>$ Br $>$ I (Ionisation energy)

(xi) F $>$ Cl $>$ Br $>$ I (Electronegativity)
\qn{5} Find the number of species / molecules having at least four identical bond lengths.

$[\mathrm{XeO_6}]^{4-}$, ClCH$_2$--CH$_2$Cl, H$_3\overset{+}{\mathrm N}-\overset{-}{\mathrm B}$H$_3$, $[\mathrm{XeF_5}]^-$, SeCl$_4$, $[\mathrm{PH_4}]^+$, AsF$_5$, $[\mathrm{SbF_6}]^-$
\qn{6} Given three conditions

(I) Involvement of $d_{z^2}$ orbital in hybridisation of central atom.

(II) All bond lengths are same

(III) Central atom is in excited state

Number of compounds that follow above two conditions but not all three conditions are:-

XeO$_4$, SF$_4$, CF$_4$, SO$_3$, XeF$_4$, XeF$_2$, BrF$_5$, IF$_7$
\sect{SECTION-III}
\qn{1} The weight (in gram) of pure potash Alum ($\mathrm{K_2SO_4.Al_2(SO_4)_3.24H_2O}$) which contains 0.64 kg oxygen. (Atomic weight of K $=39$, S $=32$, Al $=27$)

[If your answer is greater than 9 then add all the digits. Example : Ans.12 then fill $1+2=3$, Ans.123 then fill $1+2+3=6$ \& Ans. 20 then fill $2+0=2$]
\qn{2} The simplest ratio of number of oxygen atoms (O) in 16.0g of ozone (O$_3$), 28.0 g carbon monoxide (CO) and 16.0g oxygen (O$_2$) is $x:y:z$ the value of $x+y+z$ is :

(Atomic mass : C $=12$, O $=16$ and Avogadro's constant NA $=6.0\times10^{23}$ mol$^{-1}$)
\qn{3} 40 mg of gaseous substance (X$_2$) occupies 4.8 mL of volume at 1 atm and $27^\circ$C. Atomic mass of element X is : ($R=0.08$ atm L/mole-K)

(Fill the sum of digits of your answer in OMR sheet. e.g. if your answer is 1233 then fill the same as $1+2+3+3=9$)
\qn{4} Condition 1 : Shape at underlined atom is tetrahedral.

Condition 2 : Octet rule is followed for underlined atom.

Among the following the number of compounds following the above two conditions are

$\underline{\mathrm C}_2$H$_6$, $\underline{\mathrm{Si}}_2$H$_6$, Na$\underline{\mathrm B}$H$_4$, $\underline{\mathrm N}_2$O$_4$, H$_3\underline{\mathrm B}\leftarrow$N(Me)$_3$, $\underline{\mathrm B}$(OH)$_3$, $\underline{\mathrm N}_2$O$_5$, H$_2\underline{\mathrm O}_2$, Na$_2\underline{\mathrm B}_2$O$_4$(OH)$_4$
\qn{5} Count the total number of valence shell unshared electrons on I$_3^-=x$

Count the total number of p$\pi$-d$\pi$ bonds in PO$_4^{3-}=y$

Count the total number of lone pairs in XeF$_4=z$

Hence, find the value of $x-z-y$.
\qn{6} Maximum number of atom(s) in C$_3$H$_4$ which is/are in the same plane.
\columnbreak
\partsec{3}{MATHEMATICS}
\sect{SECTION-I}
\qn{1} Consider the equation $|x-6|+|x-2|-|x-4|=k$

Identify the correct statement(s) -
\begin{qoptions}
\opt{A}{Equation has exactly two positive solutions for $k=4$ whose sum is 12}
\opt{B}{Equation has exactly three positive solutions for $k=4$ whose sum is 12}
\opt{C}{Equation has exactly four solutions for $k=3$ whose sum is 17}
\opt{D}{Equation has exactly four solutions for $k=3$ whose sum is 16}
\end{qoptions}
\qn{2} If $x=\sec\phi-\tan\phi$ and $y=\operatorname{cosec}\phi+\cot\phi$, then (where $\phi\ne\dfrac{n\pi}2$, $n\in Z$)
\begin{qoptions}
\opt{A}{$x=\dfrac{y+1}{y-1}$}
\opt{B}{$x=\dfrac{y-1}{y+1}$}
\opt{C}{$y=\dfrac{1+x}{1-x}$}
\opt{D}{$xy+x-y+1=0$}
\end{qoptions}
\qn{3} If $a,b,c$ are positive numbers where $a+2b=4$ and $b+2c=5$, then
\begin{qoptions}
\opt{A}{maximum value of $ab^3$ is $\dfrac{27}8$}
\opt{B}{maximum value of $b^2c^3$ is $\dfrac{27}4$}
\opt{C}{minimum value of $\dfrac4a+\dfrac2b$ is 4}
\opt{D}{minimum value of $\dfrac9b+\dfrac2c$ is 6}
\end{qoptions}
\qn{4} If $f(n,\theta)=\sin\theta\sin^2\dfrac\theta2+\dfrac{\sin2\theta\sin^2\theta}2+\dfrac{\sin4\theta\sin^22\theta}4+\ldots+\dfrac{\sin\left(2^n\theta\right)\sin^2\left(2^{(n-1)}\theta\right)}{2^n}$; $n\in N\cup\{0\}$, then
\begin{qoptions}
\opt{A}{$f\left(5,\dfrac\pi6\right)=\dfrac{64-\sqrt3}{128}$}
\opt{B}{$f\left(5,\dfrac\pi6\right)=\dfrac{64-\sqrt3}{256}$}
\opt{C}{$f\left(2,\dfrac\pi8\right)=\dfrac{\sin(\pi/8)}2$}
\opt{D}{$f\left(2,\dfrac\pi8\right)=\sin\left(\dfrac\pi8\right)$}
\end{qoptions}
\qn{5} Let $0<A,B<\dfrac\pi2$ satisfying the equations $3\sin^2A+2\sin^2B=1$ and $3\sin2A=2\sin2B$ then which of the following is/are true -
\begin{qoptions}
\opt{A}{$A+2B=\dfrac\pi2$}
\opt{B}{$\sin A=\dfrac13$}
\opt{C}{$\cos2B=\dfrac23$}
\opt{D}{$A+2B=\dfrac\pi4$}
\end{qoptions}
\qn{6} For the equation
\par\noindent\hfill$\log_{x^2+6x+8}\left(\log_{2x^2+2x+3}\left(x^2-2x\right)\right)=0\ -$\hfill\null\par
\begin{qoptions}
\opt{A}{number of solutions is two}
\opt{B}{sum of solutions is `$-4$'}
\opt{C}{number of solution is one}
\opt{D}{sum of solutions is `$-1$'}
\end{qoptions}
\sect{SECTION-II}
\qn{1} If the minimum value of $x^2+2x+3$ is $m$ and maximum value of $-x^2+4x+6$ is $M$ then the value of $m+M$ is
\qn{2} Let $T_n=3n-2$ and $T'_m=7m+2$ be two sequences each having 2004 terms. The number of common terms to the two sequences is
\qn{3} A set of consecutive natural number from 1 to $n$ are written and one number is removed. The arithmetic mean of the remaining number is found to be $\dfrac{71}2$, then sum of all possible value(s) of the number removed from the set is $k$, then $\sqrt k$ is
\qn{4} If $\lambda_1=(0.2)^{\log_{\sqrt5}\left(\frac1{2.3.4}+\frac1{3.4.5}+\frac1{4.5.6}+\ldots..\text{upto infinite terms}\right)}$ and $\lambda_2$ is the minimum value of the expression $(3\sin x-4\cos x-10)(3\sin x+4\cos x-10)$ then $\left(\dfrac{\lambda_1-2\lambda_2}{15}\right)$ equals
\qn{5} Maximum value of
\par\noindent\hfill$\dfrac{4\sin^2x+12\sin x\cos x-12\cos^2x+3}5$\hfill\null\par
is
\qn{6} Sum of all solutions of equation $\sin\left(\pi\log_3\left(\dfrac1x\right)\right)=0$ in $[\pi,3\pi]$ is $k$ then $\dfrac k4$ is
\sect{SECTION-III}
\qn{1} If $p,q,r$ are positive numbers satisfying $pqr=3^{-5/2}$, then minimum value of $p^2+9q^2+6pq+2r^2$ is
\qn{2} \par\noindent\hfill$\dfrac{\left(\cot10^\circ-\tan10^\circ-2\operatorname{cosec}20^\circ\right)\left(\cot10^\circ+\tan10^\circ\right)}{\left(2\cot20^\circ-\cot10^\circ-\tan10^\circ\right)\operatorname{cosec}20^\circ}$\hfill\null\par
is equal to
\qn{3} Number of values of $\theta\in[0,2\pi]$ satisfying the equation $4\sin^4\theta+4\cos^4\theta=1-4x^2-4x$ for some real value of $x$, is -
\qn{4} If $y\ne0$ then number of order pair $(x,y)$ such that $x+y+\dfrac xy=\dfrac12$ and $(x+y)\dfrac xy=-\dfrac12$ is
\qn{5} Let $P(x)=4x^2+6x+7$ and $Q(y)=4y^2-12y+29$. If $P(x).Q(y)=95$, then the value of $8(x+y)$ is
\qn{6} The number of value(s) of $k$ for which equation
\par\noindent\hfill$\begin{aligned}&\left(x^2+(2k-6)x+7-3K\right)\\&\cdot\left(x^2+(2K-2)x+3K-5\right)=0\end{aligned}$\hfill\null\par
has two different pairs of equal roots is :
\end{multicols}
\end{document}
